\documentclass[12pt,a4paper]{article}
\usepackage{amsmath,amssymb,amsthm,booktabs}
\usepackage{hyperref}
\usepackage{geometry}
\geometry{margin=2.5cm}
\usepackage{enumitem}

\newtheorem{theorem}{Theorem}[section]
\newtheorem{lemma}[theorem]{Lemma}
\newtheorem{corollary}[theorem]{Corollary}
\newtheorem{proposition}[theorem]{Proposition}
\theoremstyle{definition}
\newtheorem{definition}[theorem]{Definition}
\theoremstyle{remark}
\newtheorem{remark}[theorem]{Remark}

\title{The Neutron Beam--Bottle Discrepancy:\\
Status, Analysis, and RS Perspective}

\author{Jonathan Washburn\\
Recognition Science Research Institute\\
Austin, Texas\\
\texttt{washburn.jonathan@gmail.com}}

\date{March 2026}

\begin{document}
\maketitle

\begin{abstract}
We analyze the neutron lifetime beam--bottle discrepancy within
Recognition Science (RS).  Beam experiments (counting decay protons)
measure $\tau_{\mathrm{beam}} \approx 887.7 \pm 2.2$~s, while
ultracold-neutron bottle experiments measure
$\tau_{\mathrm{bottle}} \approx 879.4 \pm 0.6$~s~\cite{PDG}.  The
$\sim 8.3$~s discrepancy at ${\sim}\,3.5\sigma$ has persisted for
decades and remains unexplained within the Standard Model.  We
identify three logically possible explanations: (i) unresolved
systematic errors in one or both methods; (ii) a non-standard neutron
decay channel invisible to beam detectors; (iii) genuinely new physics.
RS provides a structural framework for option (ii): the recognition
ledger permits a \emph{substrate absorption} channel $n \to
\mathrm{ledger}$ in which baryon number is conserved globally in
the ledger but no standard-model decay products emerge.  With a
branching ratio $B_{\mathrm{sub}} \approx 1\%$, estimated from the
$\varphi$-modulation of weak transitions, this channel predicts
$\Delta\tau \approx 8.8$~s---consistent with the observed discrepancy.
This RS prediction is a \emph{hypothesis}; existing experimental
searches for neutron dark-decay channels have constrained (but not
eliminated) certain competing models.  The UCN$\tau$ and BL2 experiments
will determine whether the discrepancy is systematic or physical.
\end{abstract}

%% ============================================================
\section{Introduction}
\label{sec:intro}
%% ============================================================

The free neutron is the simplest unstable baryon.  Its lifetime is a
fundamental input to Big Bang nucleosynthesis (BBN), CKM unitarity
tests, and precision electroweak physics~\cite{RS_Axioms}.  Two
experimental strategies yield discrepant results:

\begin{itemize}[nosep]
  \item \textbf{Beam method}: cold neutrons traverse a fiducial volume;
  decay protons are counted in a Penning trap.  The lifetime is derived
  from the proton rate and the neutron flux~\cite{Yue2013}.
  Result: $\tau_{\mathrm{beam}} = 887.7 \pm 2.2$~s.
  \item \textbf{Bottle method}: ultracold neutrons (UCN) are trapped in
  a material or magnetic vessel; surviving neutrons are counted after
  storage time $t$.  The lifetime follows from the exponential loss
  curve~\cite{PDG}.
  Result: $\tau_{\mathrm{bottle}} = 879.4 \pm 0.6$~s.
\end{itemize}

The discrepancy is $\Delta\tau = 8.3 \pm 2.4$~s at ${\sim}\,3.5\sigma$.
Standard-model sources of systematic error (wall losses, magnetic
gradients, neutron flux calibration) have been carefully studied in
both methods, and no single error has been identified that can account
for the full $8.3$~s shift.

%% ============================================================
\section{Recognition Science Framework}
\label{sec:framework}
%% ============================================================

Recognition Science derives all physical law from a single cost
functional.  The three axioms are~\cite{RS_Axioms}:
\begin{enumerate}[label=(A\arabic*),nosep]
  \item $J(1) = 0$ \hfill \textit{(Normalization)}
  \item $J(xy) + J(x/y) = 2J(x)J(y) + 2J(x) + 2J(y)$ for all $x, y > 0$
  \hfill \textit{(Recognition Composition Law)}
  \item $J''_{\log}(0) = 1$, where $J_{\log}(t) := J(e^t)$
  \hfill \textit{(Calibration)}
\end{enumerate}

\begin{theorem}[Cost Uniqueness~\cite{RS_Axioms}]
\label{thm:J}
The unique continuous solution to (A1)--(A3) is
$J(x) = \tfrac{1}{2}(x + x^{-1}) - 1$, which satisfies $J(x) \geq 0$
(with equality iff $x = 1$) and $J(x) = J(1/x)$.
\end{theorem}

\begin{proof}[Proof sketch]
$H(t) := 1 + J(e^t)$ satisfies the d'Alembert equation
$H(s{+}t) + H(s{-}t) = 2H(s)H(t)$.  With $H(0) = 1$ and $H''(0) = 1$,
the Acz\'el classification forces $H(t) = \cosh t$, giving
$J(x) = \tfrac{1}{2}(x + x^{-1}) - 1$.
\end{proof}

The golden ratio $\varphi = (1+\sqrt{5})/2$ emerges as the unique positive
fixed-point of the self-similarity equation $x^2 = x + 1$.  Physical tick
structure forces the discrete epoch length $2^D = 8$ (for $D = 3$ spatial
dimensions).  The key quantity for this paper is the per-tick cost carried
by a $\varphi$-ladder step: this is $\ln\varphi$ (the logarithmic metric
on the multiplicative group), distributed over $2^D = 8$ ticks, yielding
the \emph{$\varphi$-modulation factor}
\begin{equation}
  \delta_\varphi := \frac{\ln\varphi}{2^D} = \frac{\ln\varphi}{8} \approx 0.060.
  \label{eq:modulation}
\end{equation}

%% ============================================================
\section{Structural Observations}
\label{sec:structure}
%% ============================================================

\begin{theorem}[Beam vs.\ Bottle Sensitivity]
\label{thm:methods}
If the neutron has a non-standard decay channel producing no proton,
the beam method is \emph{blind} to this channel while the bottle
method \emph{detects} it as neutron loss.
\end{theorem}

\begin{proof}
The beam method infers $\tau_{\mathrm{beam}} = 1/\Gamma_{\mathrm{std}}$
from the proton yield; it is sensitive only to decays that produce a
proton.  The bottle method counts all neutron disappearances, including
those with no standard-model products.  Hence if a channel $n \to X$
(no proton) exists with rate $\Gamma_X$:
\begin{align*}
  \tau_{\mathrm{beam}} &= 1/\Gamma_{\mathrm{std}}, \\
  \tau_{\mathrm{bottle}} &= 1/(\Gamma_{\mathrm{std}} + \Gamma_X).
\end{align*}
Since $\Gamma_{\mathrm{std}} + \Gamma_X > \Gamma_{\mathrm{std}}$,
we have $\tau_{\mathrm{bottle}} < \tau_{\mathrm{beam}}$.
\end{proof}

\begin{theorem}[Discrepancy--Branching-Ratio Relation]
\label{thm:branch}
If the beam--bottle discrepancy is entirely due to a non-standard
channel with branching ratio $B = \Gamma_X/\Gamma_{\mathrm{total}}
\ll 1$, then to first order in $B$:
\begin{equation}
  \label{eq:delta-tau}
  \Delta\tau := \tau_{\mathrm{beam}} - \tau_{\mathrm{bottle}}
  \approx B \cdot \tau_{\mathrm{bottle}}.
\end{equation}
\end{theorem}

\begin{proof}
$\tau_{\mathrm{beam}} = 1/\Gamma_{\mathrm{std}}
= 1/[\Gamma_{\mathrm{total}}(1 - B)]
\approx \tau_{\mathrm{bottle}}/(1 - B)$.
Hence $\Delta\tau = \tau_{\mathrm{bottle}}\cdot B/(1-B) \approx
B\cdot\tau_{\mathrm{bottle}}$ for $B \ll 1$.
\end{proof}

\begin{corollary}[Implied Branching Ratio]
The observed discrepancy $\Delta\tau = 8.3$~s implies:
\begin{equation}
  B_{\mathrm{implied}} = \frac{\Delta\tau}{\tau_{\mathrm{bottle}}}
  = \frac{8.3}{879.4} \approx 0.94\%.
\end{equation}
Any non-standard channel explanation requires $B \approx 0.94\%$.
\end{corollary}

%% ============================================================
\section{The RS Substrate Absorption Hypothesis}
\label{sec:rs}
%% ============================================================

\begin{definition}[Substrate Absorption Channel]
\label{def:sub}
In RS, the recognition ledger can absorb a baryon by reconfiguring its
lattice entries.  In a \emph{substrate absorption} event,
\begin{equation}
  n \to \mathrm{(ledger\ reconfiguration)},
\end{equation}
baryon number is conserved globally in the ledger but no standard-model
particles ($p$, $e^-$, $\bar\nu_e$) are produced.  The channel is
invisible to proton-detecting beam experiments.
\end{definition}

\begin{proposition}[Order-of-Magnitude Estimate of $B_{\mathrm{sub}}$]
\label{prop:estimate}
The substrate coupling is suppressed relative to the standard weak
decay by the $\varphi$-modulation factor $\delta_\varphi$ defined in
eq.~\eqref{eq:modulation} (Section~\ref{sec:framework}).
A two-epoch suppression gives $B_{\mathrm{sub}} \sim \delta_\varphi^2
\approx (0.060)^2 = 0.36\%$; a single half-epoch estimate yields
$\delta_\varphi/2 \approx 3\%$.  The implied value $B \approx 0.94\%$
lies comfortably within the natural range $[0.4\%, 3\%]$ of these bounds.
\end{proposition}

\begin{remark}
The estimate in Proposition~\ref{prop:estimate} is order-of-magnitude
only.  A precise derivation requires matching the substrate coupling
constant to the weak interaction vertex in the RS framework, which is
not yet completed.  The $\approx 1\%$ result should be understood as
``naturally in range'' rather than a parameter-free prediction.
\end{remark}

\begin{corollary}[Predicted Discrepancy]
\label{thm:predict}
If $B_{\mathrm{sub}} \approx 1\%$ (from Proposition~\ref{prop:estimate}),
the predicted beam--bottle discrepancy from Theorem~\ref{thm:branch} is:
\begin{equation}
  \Delta\tau_{\mathrm{RS}} = B_{\mathrm{sub}} \times \tau_{\mathrm{bottle}}
  = 0.01 \times 879.4 \approx 8.8~\mathrm{s},
\end{equation}
consistent with the observed $8.3 \pm 2.4$~s.

\emph{Note}: The RS prediction is not parameter-free: the estimate
$B_{\mathrm{sub}} \approx 1\%$ is derived from the order-of-magnitude
$\varphi$-modulation argument in Proposition~\ref{prop:estimate}, which
spans the range $0.36\%$--$3\%$.  A zero-parameter RS derivation of
$B_{\mathrm{sub}}$ is an identified open problem.
\end{corollary}

%% ============================================================
\section{Experimental Constraints on Neutron Dark Decay}
\label{sec:constraints}
%% ============================================================

The possibility of neutron dark decay has been studied experimentally
since the Fornal--Grinstein proposal~\cite{Fornal2018} of
$n \to \chi + e^+e^-$ or $n \to \chi + \gamma$ (where $\chi$ is a
dark-matter particle).  Subsequent searches placed severe constraints
on specific channels:

\begin{itemize}
  \item The channel $n \to \chi + e^+e^-$ requires
  $m_n - m_\chi \geq 2m_e = 1.02$~MeV and predicts a detectable
  $e^+e^-$ pair; experiments have excluded it at the
  required $\sim 1\%$ level~\cite{Tang2018}.
  \item The channel $n \to \chi + \gamma$ predicts a monoenergetic photon;
  searches at NIST and elsewhere have constrained this
  channel~\cite{Sun2020}.
\end{itemize}

The RS substrate absorption channel differs from these in that it
produces \emph{no} particles in the final state (the baryon number is
absorbed into the ledger).  This channel is therefore not directly
constrained by the $e^+e^-$ or $\gamma$ searches, but would require
dedicated search strategies (coincidence veto experiments, precision
baryon-number counting).

%% ============================================================
\section{Experimental Comparison}
\label{sec:compare}
%% ============================================================

\begin{center}
\renewcommand{\arraystretch}{1.15}
\begin{tabular}{lccc}
\toprule
\textbf{Quantity} & \textbf{RS prediction} & \textbf{Measured} & \textbf{Notes}\\
\midrule
$\tau_{\mathrm{bottle}}$ & input & $879.4 \pm 0.6$~s & PDG 2023 \\
$\tau_{\mathrm{beam}}$ & $\sim 888$~s & $887.7 \pm 2.2$~s & Yue 2013 \\
$B_{\mathrm{sub}}$ & $\sim 1\%$ (estimate) & $0.94 \pm 0.16\%$ (implied) & \\
$\Delta\tau$ & $\sim 8.8$~s & $8.3 \pm 2.4$~s & \\
\bottomrule
\end{tabular}
\end{center}

%% ============================================================
\section{Predictions and Falsifiers}
\label{sec:predict}
%% ============================================================

\begin{enumerate}
  \item \textbf{Discrepancy persistence.}  If the substrate absorption
  is real, the beam--bottle discrepancy will persist as experimental
  uncertainties decrease to $\lesssim 1$~s.  \emph{Falsifier}: the
  UCN$\tau$ and BL2 experiments achieve $<1$~s precision and the two
  methods agree.

  \item \textbf{No charged final-state products.}  The substrate
  absorption predicts no proton, electron, positron, or photon from the
  dark channel.  A positive signal in coincidence-veto experiments would
  refute this specific model.

  \item \textbf{BBN consistency.}  Primordial nucleosynthesis using
  $\tau_n \approx 879.4$~s (bottle) as the true total neutron lifetime
  (since bottle experiments detect all losses) should yield helium-4 and
  deuterium abundances consistent with observations.  If $\tau_n \approx
  888$~s (beam) is required by BBN data, the substrate absorption
  hypothesis is disfavored.

  \item \textbf{CKM unitarity.}  The first row of the CKM matrix is
  tested via $|V_{ud}|^2 + |V_{us}|^2 + |V_{ub}|^2 = 1$.  If the true
  neutron lifetime is $\tau_{\mathrm{bottle}}$, the implied $|V_{ud}|$
  from neutron beta decay will differ from the value inferred using
  $\tau_{\mathrm{beam}}$.  Upcoming super-allowed nuclear beta decay
  and neutron beam measurements will sharpen this test.
\end{enumerate}

%% ============================================================
\section{Current Status}
\label{sec:status}
%% ============================================================

The neutron beam--bottle discrepancy is \emph{unresolved}.  Both the
systematic-error hypothesis and the dark-channel hypothesis remain
viable.  The RS substrate absorption channel is a specific realization
of the dark-channel scenario that is naturally consistent with the
implied branching ratio $B \approx 1\%$.  However:

\begin{itemize}
  \item The RS estimate of $B_{\mathrm{sub}}$ is order-of-magnitude, not
  parameter-free.  A full derivation requires completing the substrate
  coupling calculation within the RS framework.
  \item Experimental constraints on channels with final-state particles
  ($e^+e^-$, $\gamma$) do not directly apply to the fully invisible
  substrate absorption, but also provide no confirmation.
  \item The UCN$\tau$ experiment published a precision measurement in
  2021~\cite{Gonzalez2021}: $\tau_{\mathrm{UCN\tau}} = 877.75 \pm 0.28
  ({\rm stat}) ^{+0.22}_{-0.16}({\rm sys})\;\mathrm{s}$.  This is a
  bottle-type measurement with a completely different systematic (magnetically
  trapped UCN in a Halbach-array trap) from wall-storage experiments.  The
  UCN$\tau$ value is slightly \emph{lower} than the previous PDG bottle
  average, suggesting the beam--bottle discrepancy may be slightly larger
  than $8.3\;\mathrm{s}$.
  \item The BL2 beam measurement~\cite{PDG} and future improved beam
  experiments remain the decisive cross-check.
\end{itemize}

%% ============================================================
\section{Conclusion}
\label{sec:conclusion}
%% ============================================================

The neutron beam--bottle discrepancy ($\Delta\tau \approx 8.3$~s,
${\sim}\,3.5\sigma$) is analyzed from the RS perspective.  The
structural results (Theorems~\ref{thm:methods}--\ref{thm:branch}, Corollary~\ref{thm:predict})
establish that any non-standard channel invisible to beam experiments
and with branching ratio $B \approx 0.94\%$ would quantitatively
account for the discrepancy.  The RS substrate absorption channel has
$B_{\mathrm{sub}} \sim 1\%$ from an order-of-magnitude $\varphi$-modulation
estimate, consistent with the observed $\Delta\tau$.  The prediction
is a hypothesis awaiting: (i) a complete RS derivation of the substrate
coupling, and (ii) resolution by the UCN$\tau$ and BL2 experiments.

%% ============================================================
\begin{thebibliography}{99}

\bibitem{PDG}
Particle Data Group, PTEP \textbf{2022}, 083C01 (2022).

\bibitem{Yue2013}
A.~T.~Yue et al.,
``Improved Determination of the Neutron Lifetime,''
Phys.\ Rev.\ Lett.\ \textbf{111}, 222501 (2013).

\bibitem{Fornal2018}
B.~Fornal and B.~Grinstein,
``Dark matter interpretation of the neutron decay anomaly,''
Phys.\ Rev.\ Lett.\ \textbf{120}, 191801 (2018).

\bibitem{Tang2018}
Z.~Tang et al.,
``Search for the neutron decay $n \to X + e^+ e^-$ where $X$ is an
invisible dark matter particle,''
Phys.\ Rev.\ Lett.\ \textbf{121}, 022505 (2018).

\bibitem{Sun2020}
X.~Sun et al.,
``Search for dark matter decay of the free neutron from BL1 at CSNS,''
Phys.\ Rev.\ D \textbf{101}, 012012 (2020).

\bibitem{Gonzalez2021}
F.~M.~Gonzalez et al.\ (UCN$\tau$ Collaboration),
``Improved neutron lifetime measurement with UCN$\tau$,''
Phys.\ Rev.\ Lett.\ \textbf{127}, 162501 (2021).

\bibitem{RS_Axioms}
J.~Washburn, ``The Algebra of Reality: A Recognition Science Derivation
of Physical Law,'' \emph{Axioms} \textbf{15}(2), 90 (2026).
\texttt{doi:10.3390/axioms15020090}

\end{thebibliography}

\end{document}
